%% =============================================================================
%% chapters/14-configuration-hierarchy.tex
%% PART V: Configuration Architecture
%% Chapter 14: Declarative Configuration Hierarchy (Environment, ConfigMaps, Secrets, Vault)
%% =============================================================================

\chapter{Declarative Configuration Hierarchy}
\label{chap:configuration_hierarchy}

Harmonia implements a 5-stage declarative configuration hierarchy ensuring clean separation between code, infrastructure settings, and sensitive cryptographic secrets.

\section{The 5-Stage Configuration Hierarchy}

Configuration properties are resolved according to strict precedence rules:
\begin{enumerate}
    \item \textbf{Stage 1: Application Base Defaults}: Hardcoded baseline fallbacks embedded within Java \texttt{application.properties} or WildFly \texttt{standalone.xml}.
    \item \textbf{Stage 2: Declarative ConfigMaps}: Environment-specific parameters (e.g., hostnames, port bindings, timeouts, queue names) declared in Kubernetes ConfigMaps.
    \item \textbf{Stage 3: Secret Injection \& Vault}: Sensitive credentials (database passwords, Artemis security credentials, TLS private keys) mounted as Kubernetes Secrets or populated from HashiCorp Vault / Ansible Vault.
    \item \textbf{Stage 4: Container Environment Variables}: Explicit environment variables injected into container pods by Kubernetes manifest overlays.
    \item \textbf{Stage 5: Command-Line Overrides}: Ephemeral CLI flags and JVM system properties (\texttt{-Dprop=value}) used for localized diagnostic sessions.
\end{enumerate}

\section{Secret Management \& Zero-Plaintext Policy}

Harmonia strictly forbids plaintext passwords or API keys in source code or unencrypted Git repositories:
\begin{itemize}
    \item In development/testing, secrets are managed via sealed Kubernetes Secret manifests.
    \item In production, secrets are dynamically fetched at startup using external secret operators or Ansible Vault injection during deployment.
\end{itemize}

\newpage
